\section{Interception and deadline-fuel games}
\label{sec:interception}
This section restores position-only interception, needed for contact or
salvage races. Boarding instead requires velocity matching. All closed
forms here use the ideal engine and instantaneous nonanticipative response;
they are not finite-band game values. Write $r=x_E-x_P$, $w=v_E-v_P$,
and take capture to mean $|r|\le R$. Resources in this section may be
physical, with $g=-\mu x/|x|^3$, or consistently normalized.

\subsection{The four problems and verification}
Let $z=(x_P,v_P,x_E,v_E,b_P,b_E)$ and
\[
 \mathcal D U=U_{x_P}\cdot v_P+U_{v_P}\cdot g(x_P)
              +U_{x_E}\cdot v_E+U_{v_E}\cdot g(x_E).
\]
The Isaacs operator for pursuer running cost $\ell|a_P|^2$ is
\begin{equation}
\begin{split}
\mathcal H_\ell U={}&\mathcal D U+
 \inf_{a_P}\{U_{v_P}\cdot a_P+(\ell-U_{b_P})|a_P|^2\}\\
 &+\sup_{a_E}\{U_{v_E}\cdot a_E-U_{b_E}|a_E|^2\}.
\end{split}
\end{equation}
At a zero resource face the corresponding control is zero. Where
$c_P=\ell-U_{b_P}>0$ and $c_E=U_{b_E}>0$, this becomes
\begin{equation}
 \begin{gathered}
 \mathcal H_\ell U=\mathcal D U-
 \frac{|U_{v_P}|^2}{4c_P}+\frac{|U_{v_E}|^2}{4c_E},\\
 a_P^*=-\frac{U_{v_P}}{2c_P},\qquad a_E^*=\frac{U_{v_E}}{2c_E}.
 \end{gathered}
\end{equation}
Zero or wrong-sign coefficients require evaluating the original extrema,
including infinite values, rather than dividing by zero.
Problems A1 and B1 minimize guaranteed interception time in free space
and Kepler gravity, respectively: $1+\mathcal H_0V=0$, $V=0$ on capture.
Problems A2 and B2 minimize worst-case pursuer effort to capture within
remaining time $s$: $W_s=\mathcal H_1W$, with $W=0$ on capture and
$W(0,z)=+\infty$ outside capture. Failure to capture by the deadline costs
$+\infty$; the budget state remains a hard constraint.

\begin{proposition}[Classical verification]
Suppose a nonnegative differentiable solution satisfies the stated PDE,
target and resource boundary conditions. Suppose its selectors generate
well-posed admissible responses throughout the relevant domain. For time,
the pursuer selector guarantees capture and the evader selector gives the
opposing time bound. For deadline effort, assume additionally that the
pursuer selector guarantees capture before the deadline. Then the solution
is the respective value, provided the boundary limits used at capture
and exhaustion exist.
\end{proposition}
Indeed, along the time pursuer selector, $dV/dt\le-1$ against every
evader, whereas along the evader selector $dV/dt\ge-1$ against every
pursuer. Integrate to capture; nonnegativity supplies the time upper
bound. For remaining time $s'= -1$, along the deadline pursuer selector,
$dW/dt+|a_P|^2\le0$, and along the evader selector the inequality reverses.
Integration gives the opposing effort bounds; noncapture has infinite
cost. Localization and boundary limits are essential. Nonsmooth values
require viscosity and comparison hypotheses \cite{EvansSouganidis1984};
this verification statement does not assert existence for an unrestricted
singular orbital domain.

\subsection{Closed forms and distinctions}
Against a ballistic evader in free space, set
$d(h)=(1-R/|r+hw|)_+(r+hw)$, with $d=0$ at zero separation. The exact
prescribed-time interception effort and control are
\begin{equation}
 c(h)=\frac{3|d(h)|^2}{h^3},\qquad a_P(t)=\frac{3(h-t)}{h^3}d(h).
\end{equation}
These follow by projecting the required position onto the capture ball
and using the single moment $\int_0^h(h-t)a_Pdt=d$. Deadline fuel is
$\inf_{0<h\le T}c(h)$ and earliest interception is
$\inf\{h:c(h)\le b_P\}$. In particular, if $c(h)>b_P$ for every
$0<h\le T$, coasting is a zero-fuel evader win, even when fuel is available.

For active interception with $w=0$, $D=|r|-R>0$, $b_P>b_E$ and
$b_P\ge4b_E$, the exact guaranteed time is
\begin{equation}
 H_I=\left[\frac{3D^2}{(\sqrt{b_P}-\sqrt{b_E})^2}\right]^{1/3}.
\end{equation}
Put $k=\sqrt{b_E}/(\sqrt{b_P}-\sqrt{b_E})\le1$ and
$a_0(t)=3D(1-t/H_I)r/(H_I^2|r|)$. The response $a_P=a_E+a_0$
gives the upper bound by Minkowski. The opponent $a_E=ka_0$ prevents an
earlier hit: at $t=sH_I< H_I$ any pursuer needs at least
$C s^{-3}[1+kf(s)]^2$, where $C=3D^2/H_I^3$ and
$f(s)=(3s^2-s^3)/2$. This exceeds $(1+k)^2C=b_P$ because
\[
\begin{split}
1+kf(s)-(1+k)s^{3/2}
={}&(1-k)(1-s^{3/2})\\
 &+k[1+f(s)-2s^{3/2}]>0;
\end{split}
\]
use $f(s)\ge s^3$. For point capture ($R=0$), full-state boarding takes
$4^{1/3}H_I$ under these same hypotheses.

For a prescribed exact endpoint at time $h$, let
\begin{equation}
 \begin{gathered}
 C_I=3|r+hw|^2/h^3,\\
 C_J=12|r|^2/h^3+12r\cdot w/h^2+4|w|^2/h.
 \end{gathered}
\end{equation}
The endpoint minimax pursuer effort is
\begin{equation}
 M_{I,J}(h)=(\sqrt{b_E}+\sqrt{C_{I,J}(h)})^2.
\end{equation}
The copying response gives the upper bound. Choosing the evader control
proportional to the minimum-norm endpoint correction gives the lower
bound by orthogonal projection onto the endpoint moment constraints.
This is an exact-time endpoint result. Minimizing it over deadlines gives
a sufficient first-hit bound, not automatically the first-hit value.
For perpetual avoidance of exact rendezvous at a nonzero phase state,
$b_E\ge b_P$ gives least worst-case evader expenditure
$F^J_\infty=b_P$. The parity feedback in Appendix~\ref{app:game} attains
the upper bound; a smaller uniform expenditure cap would give the
pursuer a strictly positive resource gap and finite-time rendezvous.

\paragraph{Nonattainment matters.}
Strict position-only escape can have an infimum without a minimizing
control. An exhausted stationary pursuer at the origin and an evader
coasting from $(-L,R)$ with velocity $(u,0)$ meet the capture boundary
at $L/u$. Any arbitrarily small positive transverse thrust before that
time avoids it forever, at cost tending to zero; zero thrust loses.
There are also positive unattained infima. In one dimension take point
capture, separation $L>0$, closing velocity $-u<0$, and an exhausted
pursuer. The critical effort is $4u^3/(9L)$: the endpoint lower bound
$3(ut-L)^2/t^3$ is largest at $t=3L/u$. Equality forces contact there.
The safe profiles $r(t)=\delta+(L-\delta)(1-t/t_\delta)^3$ for
$0\le t\le t_\delta=3(L-\delta)/u$, followed by rest at $\delta>0$,
approach that cost as $\delta\downarrow0$ without attaining it.

\subsection{Five-state reduction and scaling}
In isotropic free space a value depends only on
$(\rho_G,\xi_G,\omega_G,b_P,b_E)=(|r|^2,r\cdot w,|w|^2,b_P,b_E)$,
where $\xi_G^2\le\rho_G\omega_G$. The drift is
$2\xi_G U_{\rho_G}+\omega_G U_{\xi_G}$ and
\begin{equation}
 |U_w|^2=\rho_G U_{\xi_G}^2+4\xi_G U_{\xi_G}U_{\omega_G}
                      +4\omega_G U_{\omega_G}^2.
\end{equation}
Both controlled velocity gradients have this squared norm. These chain
rules reduce the free-space Isaacs equation to five states in $d=2$ or
$3$. Under $t\mapsto\alpha t$, $r\mapsto\beta r$,
$w\mapsto(\beta/\alpha)w$, resources scale as $\beta^2/\alpha^3$.
Time values scale by $\alpha$, and deadline-fuel values by
$\beta^2/\alpha^3$, with radius and deadline scaled as well. Scaling
every strategy and opponent proves the rule for values, not only paths.

Finally, restore dimensional mean motion $n$ in the CW matrices. As
$n\to0$, $\Phi_n$ and $W_n$ converge to the double-integrator matrices.
Theorem~1 therefore reduces to the free-space cubic
\begin{equation}
 (\sqrt{b_P}-\sqrt{b_E})^2H^3-4|w|^2H^2
       -12r\cdot w H-12|r|^2=0.
\end{equation}
One takes this physical limit before setting the normalized mean motion
to one.

